\documentclass[11pt]{article}
\usepackage[utf8]{inputenc}
\usepackage[T1]{fontenc}
\usepackage[spanish,mexico]{babel}
\usepackage{amsmath,amssymb}
\usepackage{geometry}
\usepackage{enumitem}
\usepackage{multicol}
\usepackage{fancyhdr}
\geometry{letterpaper,margin=1.7cm}
\setlength{\parindent}{0pt}
\pagestyle{fancy}
\fancyhf{}
\lhead{Álgebra Lineal}
\rhead{Fluidez Aritmética --- Versión B}
\cfoot{\thepage}

\begin{document}
\begin{center}
{\Large\bfseries Prueba de Fluidez Aritmética para Álgebra}\\[1mm]
{\large Versión B --- Diagnóstico preliminar}
\end{center}

\textbf{Nombre:} \rule{9cm}{0.4pt}\hfill
\textbf{Fecha:} \rule{3cm}{0.4pt}

\medskip
\textbf{Instrucciones.} Resuelva mentalmente cada reactivo y escriba únicamente
el resultado. No utilice calculadora. Simplifique toda fracción hasta obtener
una fracción irreducible. En los reactivos de MCD escriba el máximo común divisor.

\medskip
\textbf{Nota.} El instrumento está diseñado como diagnóstico de clase; no es
una prueba psicométrica estandarizada.

\subsection*{I. Cálculo aritmético inmediato}
\begin{multicols}{2}\begin{enumerate}[leftmargin=*,start=1]
\item $\displaystyle 8+7=$ \rule{2.0cm}{0.4pt}
\item $\displaystyle 14-9=$ \rule{2.0cm}{0.4pt}
\item $\displaystyle 6\cdot7=$ \rule{2.0cm}{0.4pt}
\item $\displaystyle 36\div4=$ \rule{2.0cm}{0.4pt}
\item $\displaystyle 9+6-4=$ \rule{2.0cm}{0.4pt}
\item $\displaystyle 5\cdot8-7=$ \rule{2.0cm}{0.4pt}
\item $\displaystyle 18\div3+4=$ \rule{2.0cm}{0.4pt}
\item $\displaystyle 2(7-3)=$ \rule{2.0cm}{0.4pt}
\item $\displaystyle 24\div(5-2)=$ \rule{2.0cm}{0.4pt}
\item $\displaystyle 3^2+5=$ \rule{2.0cm}{0.4pt}
\item $\displaystyle 17-(-4)=$ \rule{2.0cm}{0.4pt}
\item $\displaystyle -6+11=$ \rule{2.0cm}{0.4pt}
\item $\displaystyle (-3)(5)=$ \rule{2.0cm}{0.4pt}
\item $\displaystyle (-4)(-6)=$ \rule{2.0cm}{0.4pt}
\item $\displaystyle 21-3(5)=$ \rule{2.0cm}{0.4pt}
\end{enumerate}\end{multicols}
\subsection*{II. División mental por $3$, $5$ y $7$}
\begin{multicols}{2}\begin{enumerate}[leftmargin=*,start=16]
\item $\displaystyle 18\div3=$ \rule{2.0cm}{0.4pt}
\item $\displaystyle 27\div3=$ \rule{2.0cm}{0.4pt}
\item $\displaystyle 33\div3=$ \rule{2.0cm}{0.4pt}
\item $\displaystyle 42\div3=$ \rule{2.0cm}{0.4pt}
\item $\displaystyle 51\div3=$ \rule{2.0cm}{0.4pt}
\item $\displaystyle 25\div5=$ \rule{2.0cm}{0.4pt}
\item $\displaystyle 35\div5=$ \rule{2.0cm}{0.4pt}
\item $\displaystyle 45\div5=$ \rule{2.0cm}{0.4pt}
\item $\displaystyle 55\div5=$ \rule{2.0cm}{0.4pt}
\item $\displaystyle 65\div5=$ \rule{2.0cm}{0.4pt}
\item $\displaystyle 21\div7=$ \rule{2.0cm}{0.4pt}
\item $\displaystyle 35\div7=$ \rule{2.0cm}{0.4pt}
\item $\displaystyle 42\div7=$ \rule{2.0cm}{0.4pt}
\item $\displaystyle 49\div7=$ \rule{2.0cm}{0.4pt}
\item $\displaystyle 63\div7=$ \rule{2.0cm}{0.4pt}
\item $\displaystyle 54\div3=$ \rule{2.0cm}{0.4pt}
\item $\displaystyle 75\div5=$ \rule{2.0cm}{0.4pt}
\item $\displaystyle 77\div7=$ \rule{2.0cm}{0.4pt}
\item $\displaystyle 84\div7=$ \rule{2.0cm}{0.4pt}
\item $\displaystyle 105\div7=$ \rule{2.0cm}{0.4pt}
\end{enumerate}\end{multicols}
\subsection*{III. Máximo común divisor}
\begin{multicols}{2}\begin{enumerate}[leftmargin=*,start=36]
\item $\displaystyle \operatorname{MCD}(12,18)=$ \rule{2.0cm}{0.4pt}
\item $\displaystyle \operatorname{MCD}(15,25)=$ \rule{2.0cm}{0.4pt}
\item $\displaystyle \operatorname{MCD}(14,21)=$ \rule{2.0cm}{0.4pt}
\item $\displaystyle \operatorname{MCD}(18,27)=$ \rule{2.0cm}{0.4pt}
\item $\displaystyle \operatorname{MCD}(20,30)=$ \rule{2.0cm}{0.4pt}
\item $\displaystyle \operatorname{MCD}(21,35)=$ \rule{2.0cm}{0.4pt}
\item $\displaystyle \operatorname{MCD}(24,36)=$ \rule{2.0cm}{0.4pt}
\item $\displaystyle \operatorname{MCD}(28,42)=$ \rule{2.0cm}{0.4pt}
\item $\displaystyle \operatorname{MCD}(30,45)=$ \rule{2.0cm}{0.4pt}
\item $\displaystyle \operatorname{MCD}(42,63)=$ \rule{2.0cm}{0.4pt}
\item $\displaystyle \operatorname{MCD}(16,24)=$ \rule{2.0cm}{0.4pt}
\item $\displaystyle \operatorname{MCD}(25,40)=$ \rule{2.0cm}{0.4pt}
\item $\displaystyle \operatorname{MCD}(27,36)=$ \rule{2.0cm}{0.4pt}
\item $\displaystyle \operatorname{MCD}(35,50)=$ \rule{2.0cm}{0.4pt}
\item $\displaystyle \operatorname{MCD}(49,63)=$ \rule{2.0cm}{0.4pt}
\end{enumerate}\end{multicols}
\subsection*{IV. Simplificación de fracciones}
\begin{multicols}{2}\begin{enumerate}[leftmargin=*,start=51]
\item $\displaystyle \frac{14}{21}=$ \rule{2.0cm}{0.4pt}
\item $\displaystyle \frac{18}{27}=$ \rule{2.0cm}{0.4pt}
\item $\displaystyle \frac{25}{35}=$ \rule{2.0cm}{0.4pt}
\item $\displaystyle \frac{28}{42}=$ \rule{2.0cm}{0.4pt}
\item $\displaystyle \frac{30}{45}=$ \rule{2.0cm}{0.4pt}
\item $\displaystyle \frac{35}{49}=$ \rule{2.0cm}{0.4pt}
\item $\displaystyle \frac{42}{63}=$ \rule{2.0cm}{0.4pt}
\item $\displaystyle \frac{45}{60}=$ \rule{2.0cm}{0.4pt}
\item $\displaystyle \frac{50}{70}=$ \rule{2.0cm}{0.4pt}
\item $\displaystyle \frac{54}{72}=$ \rule{2.0cm}{0.4pt}
\item $\displaystyle \frac{21}{27}=$ \rule{2.0cm}{0.4pt}
\item $\displaystyle \frac{25}{30}=$ \rule{2.0cm}{0.4pt}
\item $\displaystyle \frac{35}{45}=$ \rule{2.0cm}{0.4pt}
\item $\displaystyle \frac{42}{55}=$ \rule{2.0cm}{0.4pt}
\item $\displaystyle \frac{49}{63}=$ \rule{2.0cm}{0.4pt}
\end{enumerate}\end{multicols}
\subsection*{V. Operaciones con fracciones}
\begin{multicols}{2}\begin{enumerate}[leftmargin=*,start=66]
\item $\displaystyle \frac12+\frac13=$ \rule{2.0cm}{0.4pt}
\item $\displaystyle \frac34-\frac12=$ \rule{2.0cm}{0.4pt}
\item $\displaystyle \frac23+\frac16=$ \rule{2.0cm}{0.4pt}
\item $\displaystyle \frac56-\frac13=$ \rule{2.0cm}{0.4pt}
\item $\displaystyle \frac35+\frac1{10}=$ \rule{2.0cm}{0.4pt}
\item $\displaystyle \frac78-\frac38=$ \rule{2.0cm}{0.4pt}
\item $\displaystyle \frac23+\frac14=$ \rule{2.0cm}{0.4pt}
\item $\displaystyle \frac56-\frac14=$ \rule{2.0cm}{0.4pt}
\item $\displaystyle \frac37+\frac27=$ \rule{2.0cm}{0.4pt}
\item $\displaystyle \frac59-\frac29=$ \rule{2.0cm}{0.4pt}
\item $\displaystyle \frac23\cdot\frac35=$ \rule{2.0cm}{0.4pt}
\item $\displaystyle \frac57\cdot\frac{14}{15}=$ \rule{2.0cm}{0.4pt}
\item $\displaystyle \frac34\div\frac12=$ \rule{2.0cm}{0.4pt}
\item $\displaystyle \frac67\div\frac9{14}=$ \rule{2.0cm}{0.4pt}
\item $\displaystyle \frac{14}{15}\cdot\frac5{21}=$ \rule{2.0cm}{0.4pt}
\item $\displaystyle 1-\frac38=$ \rule{2.0cm}{0.4pt}
\item $\displaystyle 2-\frac56=$ \rule{2.0cm}{0.4pt}
\item $\displaystyle \frac12+\frac13+\frac16=$ \rule{2.0cm}{0.4pt}
\item $\displaystyle \left(\frac34-\frac14\right)\div\frac12=$ \rule{2.0cm}{0.4pt}
\item $\displaystyle \left(\frac23+\frac13\right)\frac35=$ \rule{2.0cm}{0.4pt}
\end{enumerate}\end{multicols}
\subsection*{VI. Operaciones mixtas y simplificación}
\begin{multicols}{2}\begin{enumerate}[leftmargin=*,start=86]
\item $\displaystyle \frac{18}{24}+\frac14=$ \rule{2.0cm}{0.4pt}
\item $\displaystyle \frac{21}{35}+\frac25=$ \rule{2.0cm}{0.4pt}
\item $\displaystyle \frac{28}{42}-\frac16=$ \rule{2.0cm}{0.4pt}
\item $\displaystyle \frac35+\frac{15}{25}=$ \rule{2.0cm}{0.4pt}
\item $\displaystyle \frac79-\frac{14}{27}=$ \rule{2.0cm}{0.4pt}
\item $\displaystyle \frac47\cdot\frac{21}{8}=$ \rule{2.0cm}{0.4pt}
\item $\displaystyle \frac56\div\frac{10}{9}=$ \rule{2.0cm}{0.4pt}
\item $\displaystyle 1-\left(\frac23-\frac16\right)=$ \rule{2.0cm}{0.4pt}
\item $\displaystyle \left(\frac12+\frac14\right)\div\frac34=$ \rule{2.0cm}{0.4pt}
\item $\displaystyle \frac35\cdot\frac{10}{9}+\frac13=$ \rule{2.0cm}{0.4pt}
\item $\displaystyle \left(\frac78-\frac14\right)\div\frac58=$ \rule{2.0cm}{0.4pt}
\item $\displaystyle \frac{15}{21}+\frac27=$ \rule{2.0cm}{0.4pt}
\item $\displaystyle \frac{18}{30}\cdot\frac53=$ \rule{2.0cm}{0.4pt}
\item $\displaystyle \frac{21}{28}\div\frac38=$ \rule{2.0cm}{0.4pt}
\item $\displaystyle 2-\frac35-\frac25=$ \rule{2.0cm}{0.4pt}
\end{enumerate}\end{multicols}

\newpage
\section*{Clave de respuestas --- uso del profesor}
\subsection*{I. Cálculo aritmético inmediato}
\begin{multicols}{4}\begin{enumerate}[leftmargin=*,start=1]
\item $\displaystyle 15$
\item $\displaystyle 5$
\item $\displaystyle 42$
\item $\displaystyle 9$
\item $\displaystyle 11$
\item $\displaystyle 33$
\item $\displaystyle 10$
\item $\displaystyle 8$
\item $\displaystyle 8$
\item $\displaystyle 14$
\item $\displaystyle 21$
\item $\displaystyle 5$
\item $\displaystyle -15$
\item $\displaystyle 24$
\item $\displaystyle 6$
\end{enumerate}\end{multicols}
\subsection*{II. División mental por $3$, $5$ y $7$}
\begin{multicols}{4}\begin{enumerate}[leftmargin=*,start=16]
\item $\displaystyle 6$
\item $\displaystyle 9$
\item $\displaystyle 11$
\item $\displaystyle 14$
\item $\displaystyle 17$
\item $\displaystyle 5$
\item $\displaystyle 7$
\item $\displaystyle 9$
\item $\displaystyle 11$
\item $\displaystyle 13$
\item $\displaystyle 3$
\item $\displaystyle 5$
\item $\displaystyle 6$
\item $\displaystyle 7$
\item $\displaystyle 9$
\item $\displaystyle 18$
\item $\displaystyle 15$
\item $\displaystyle 11$
\item $\displaystyle 12$
\item $\displaystyle 15$
\end{enumerate}\end{multicols}
\subsection*{III. Máximo común divisor}
\begin{multicols}{4}\begin{enumerate}[leftmargin=*,start=36]
\item $\displaystyle 6$
\item $\displaystyle 5$
\item $\displaystyle 7$
\item $\displaystyle 9$
\item $\displaystyle 10$
\item $\displaystyle 7$
\item $\displaystyle 12$
\item $\displaystyle 14$
\item $\displaystyle 15$
\item $\displaystyle 21$
\item $\displaystyle 8$
\item $\displaystyle 5$
\item $\displaystyle 9$
\item $\displaystyle 5$
\item $\displaystyle 7$
\end{enumerate}\end{multicols}
\subsection*{IV. Simplificación de fracciones}
\begin{multicols}{4}\begin{enumerate}[leftmargin=*,start=51]
\item $\displaystyle \frac23$
\item $\displaystyle \frac23$
\item $\displaystyle \frac57$
\item $\displaystyle \frac23$
\item $\displaystyle \frac23$
\item $\displaystyle \frac57$
\item $\displaystyle \frac23$
\item $\displaystyle \frac34$
\item $\displaystyle \frac57$
\item $\displaystyle \frac34$
\item $\displaystyle \frac79$
\item $\displaystyle \frac56$
\item $\displaystyle \frac79$
\item $\displaystyle \frac{42}{55}$
\item $\displaystyle \frac79$
\end{enumerate}\end{multicols}
\subsection*{V. Operaciones con fracciones}
\begin{multicols}{4}\begin{enumerate}[leftmargin=*,start=66]
\item $\displaystyle \frac56$
\item $\displaystyle \frac14$
\item $\displaystyle \frac56$
\item $\displaystyle \frac12$
\item $\displaystyle \frac7{10}$
\item $\displaystyle \frac12$
\item $\displaystyle \frac{11}{12}$
\item $\displaystyle \frac7{12}$
\item $\displaystyle \frac57$
\item $\displaystyle \frac13$
\item $\displaystyle \frac25$
\item $\displaystyle \frac23$
\item $\displaystyle \frac32$
\item $\displaystyle \frac43$
\item $\displaystyle \frac29$
\item $\displaystyle \frac58$
\item $\displaystyle \frac76$
\item $\displaystyle 1$
\item $\displaystyle 1$
\item $\displaystyle \frac35$
\end{enumerate}\end{multicols}
\subsection*{VI. Operaciones mixtas y simplificación}
\begin{multicols}{4}\begin{enumerate}[leftmargin=*,start=86]
\item $\displaystyle 1$
\item $\displaystyle 1$
\item $\displaystyle \frac12$
\item $\displaystyle \frac65$
\item $\displaystyle \frac7{27}$
\item $\displaystyle \frac32$
\item $\displaystyle \frac34$
\item $\displaystyle \frac12$
\item $\displaystyle 1$
\item $\displaystyle 1$
\item $\displaystyle 1$
\item $\displaystyle 1$
\item $\displaystyle 1$
\item $\displaystyle 2$
\item $\displaystyle 1$
\end{enumerate}\end{multicols}
\end{document}
